% !PRINTED 146
\weberpage{146}
%<<CONT
und dieselben Vertauschungen auch auf $\mathrm{sn}(u \pm v)$, $\mathrm{cn}(u \pm v)$, $\mathrm{dn}(u \pm v)$ anwendet. So ergeben sich z.\,B. die zw\"olf Formeln (1) bis (12) aus den drei Formeln (16).

Wendet man diese Vertauschungen auf die Relationen (18), (19), (20) selber an, so ergibt sich:
\weq{22}{\begin{aligned}
\mathrm{sn}(u+2K) &= -\mathrm{sn}\,u, & \mathrm{sn}(u+2iK') &= \mathrm{sn}\,u,\\
\mathrm{cn}(u+2K) &= -\mathrm{cn}\,u, & \mathrm{cn}(u+2iK') &= -\mathrm{cn}\,u,\\
\mathrm{dn}(u+2K) &= \mathrm{dn}\,u, & \mathrm{dn}(u+2iK') &= -\mathrm{dn}\,u.
\end{aligned}}
Die gemeinschaftlichen Perioden dieser drei Funktionen sind also $4K$, $4iK'$; au\ss erdem hat aber $\mathrm{sn}\,u$ die Periode $2iK'$, $\mathrm{dn}\,u$ die Periode $2K$, $\mathrm{cn}\,u$ die Periode $2K + 2iK'$.

Wir geben nur dem Satz \"uber $\Theta$-Funktionen (\S\,21) einen anderen Ausdruck, indem wir den folgenden Satz aussprechen:

\gesperrt{Jede doppelt periodische Funktion von $v$ mit den Perioden $2K$, $2iK'$ (die \"uberall den Charakter einer rationalen Funktion hat) ist, wenn sie eine gerade Funktion ist, eine rationale Funktion von $\mathrm{sn}^2v$, und wenn sie eine ungerade Funktion ist, das Produkt von $\mathrm{sn}\,v\,\mathrm{cn}\,v\,\mathrm{dn}\,v$ mit einer rationalen Funktion von $\mathrm{sn}^2v$.}

So kann man jeden Satz, der sich auf homogene Funktionen nullter Ordnung von $\vartheta$-Funktionen bezieht, in einen Satz \"uber elliptische Funktionen verwandeln, und erh\"alt daraus wieder entsprechende S\"atze \"uber elliptische Integrale. Man erkennt sofort, da\ss{} die Additionsformeln (16) keine anderen sind, als die Formeln (17) des \S\,13 im ersten Abschnitt. In derselben Weise liefert uns die Transformationstheorie der $\vartheta$-Funktionen eine L\"osung des \gesperrt{Jacobischen} Transformationsproblems der elliptischen Integrale, wie wir es im \S\,8 dargelegt haben.

Wir wollen dies an den beiden Haupttransformationen zweiter Ordnung nachweisen, die wir im \S\,32 ebenso wie im \S\,9 als \gesperrt{Gausssche} und \gesperrt{Landensche} Transformation bezeichnet haben. Nach \S\,32 ist f\"ur die \gesperrt{Gausssche} Transformation
\weq{23}{\frac{\vartheta_{10}\left(0,\dfrac{\omega}{2}\right)\vartheta_{11}\left(u,\dfrac{\omega}{2}\right)}{\vartheta_{00}\left(0,\dfrac{\omega}{2}\right)\vartheta_{01}\left(u,\dfrac{\omega}{2}\right)} = \frac{2\,\vartheta_{01}(u)\,\vartheta_{11}(u)}{\vartheta_{01}^2(u) + \vartheta_{11}^2(u)},}
\weq{24}{\frac{\vartheta_{10}\left(0,\dfrac{\omega}{2}\right)^2}{\vartheta_{00}\left(0,\dfrac{\omega}{2}\right)^2} = \frac{2\,\vartheta_{10}\,\vartheta_{00}}{\vartheta_{00}^2 + \vartheta_{10}^2} = \frac{2\sqrt{\varkappa}}{1+\varkappa},}
%CONT>>
